%!TEX Root = ../main.tex

%%%%%%%%
\section{\label{sec:intro}Introduction and Summary}
%%%%%%%%
Every four-dimensional $\mathcal{N} = 2$ superconformal field theory (SCFT) gives rise to an associated vertex operator algebra (VOA), via the SCFT/VOA correspondence \cite{Beem:2013sza}. These vertex algebras possess a number of nice properties that make them better behaved than generic VOAs. For example, they are conjecturally quasi-lisse so their vacuum characters satisfy finite-order modular linear differential equations \cite{Beem:2017ooy,Arakawa:2016hkg} and in many cases they admit extremely nice free field realisations \cite{Beem:2019tfp,Beem:2019snk,Beem:2021jnm,Beem:2024fom}. A salient feature is that when the four-dimensional theory in question is unitary, the associated VOA inherits a novel form of unitarity, christened \emph{graded unitarity} in \cite{ArabiArdehali:2025fad}. The implications of this structure for Virasoro, $\mathfrak{sl}_2, \mathfrak{sl}_3, \mathfrak{sl}_4$ affine vertex algebras were also investigated \emph{loc. cit.}, while the more technically involved case of Zamolodchikov $\mathcal{W}_3$ algebras was examined in \cite{Beem:2026lkq}. The variations of this structure, along with a more mathematically precise formulation were elaborated in \cite{Beem:2025guj}. Some preliminary consequences of four-dimensional unitarity for the associated vertex operator algebra were studied in \cite{Liendo:2015ofa,Beem:2018duj}, which in retrospect are consequences of graded unitarity. More recently, it was established that graded unitarity of $\mathfrak{sl}_2$ at boundary admissible level $k = -2 + \frac{2}{q}$ implies graded unitarity of the $(2,q)$ Virasoro minimal models \cite{Popov:2026qxn}. 

A graded unitarity structure for a vertex algebra $\mathbb{V}$ amounts, in particular, to a conjugation operation $\rho$ (generally of order four) together with a nondegenerate half-integral good filtration $\mathfrak{R}_{\bullet}$ on $\mathbb{V}$. With respect to these ingredients, graded unitarity amounts to the requirement that a suitably defined Hermitian form is positive definite; equivalently, that the constants appearing in the two-point functions of operators with their conjugates have definite, $R$-charge-dependent signs. It is an important problem to characterise to what extent the existence of a graded unitary structure constrains the underlying vertex algebra. A direct attack on the problem (for a fixed underlying vertex algebra) is complicated by the dependence of the positivity property on the choice of $\mathfrak{R}$-filtration. At present, there is no intrinsic characterisation of the physically relevant $\mathfrak{R}$-filtration without reference to four-dimensional input. One is thus led to an interesting (but seemingly difficult) problem of constraining a choice of good filtrations subject to highly specific sign requirements for two-point functions.

If, however, an $\mathfrak{R}$-filtration is given or assumed, then graded unitarity makes extremely precise predictions about the positivity/negativity conditions for two-point functions in the underlying vertex algebra, and these are comparatively accessible. This approach was pursued in \cite{ArabiArdehali:2025fad, Beem:2026lkq} for the cases where the underlying vertex algebra is Virasoro, an affine Kac--Moody VOA of type $\mathfrak{sl}_n$ with $n=2,3,4$, or the Zamolodchikov $\mathcal{W}_3$ VOA. In these works, a simple but physically motivated choice of $\mfR$-filtration was assumed and consequently all but a discrete set of values for the central charge/level were ruled out by imposing the relevant positivity bounds level by level. These stringent restrictions on the central charge/level were obtained by examining the positivity constraints visible at the level of the Kac determinant at a given conformal level (and $\mathfrak{g}$-weight in the Kac--Moody case). While the full set of positivity constraints from graded unitarity are much richer, it is particularly satisfying that determinant-based constraints provide such strict restrictions.

The purpose of the present paper is to generalise such a classification to \emph{all} simply-laced affine Kac--Moody algebras. There are several families of four-dimensional theories which are known, or conjectured, to give rise to simply-laced affine Kac--Moody algebras. In particular, the VOAs associated with the Kodaira/Deligne--Cvitanovi\'{c} rank-one SCFTs \cite{Seiberg:1994aj, Argyres:1995xn, Minahan:1996fg, Minahan:1996cj} are given by affine Kac--Moody VOAs $V_k(\gf)$ at level $k = -\frac{h^\vee_\gf}{6}-1$, where $h^\vee$ is the dual Coxeter number of the corresponding simple Lie algebra $\gf$. For simply-laced $\gf$, these include the family,
\begin{equation}
    \mathfrak{a}_1 \subset \mathfrak{a}_2 \subset \mathfrak{d}_4 \subset \mathfrak{e}_6 \subset \mathfrak{e}_7 \subset \mathfrak{e}_8~,
\end{equation}
so the respective affine Kac--Moody VOAs are expected to be graded unitary. 

Furthermore, a plethora of generalised Argyres--Douglas (AD) theories \cite{Argyres:1995jj, Argyres:1995xn, Xie:2012hs} conjecturally give rise to affine Kac--Moody VOAs. Specifically, the theories of interest are associated to irregular Hitchin systems of type $ADE$ with a full regular puncture and an irregular puncture \cite{Wang:2015mra}. We denote these theories by $D^b_q(\gf)$, where $b$ and $q$ parametrise the irregular singularity, and $\gf$ is simply laced.\footnote{Here, we use a different parametrisation from the one used in \cite{Wang:2015mra,Xie:2016evu, Song:2017oew}, where the theories are parametrised by $k$ and $b$. The two conventions are related by $k = -b+q$.} Largely based on the requirement that the four-dimensional theory has precisely $\gf$-flavour symmetry, and the satisfaction of the Sugawara relation for the central charge, the associated VOAs for the $D_q^{b}(\gf)$ generalised Argyres--Douglas SCFTs without extra mass parameters/flavour symmetries from the irregular puncture, are conjectured to be the affine Kac--Moody VOAs \cite{Xie:2016evu, Song:2017oew, Xie:2017aqx},
\begin{equation} \label{eq:DqbgVOA}
    \mathbb{V}[D_q^{b}(\gf)] = V_{k}(\gf), \qquad k = -h^\vee + \frac{b}{q}~. 
\end{equation}
It is worth emphasising that (non-trivial) finite extensions by modules of affine Kac--Moody algebras do \emph{not} affect the central charge $c_{4d}$ of the putative four-dimensional theory. Consequently, the agreement of the central charge alone does not distinguish the conjectured identification in \eqref{eq:DqbgVOA} from possible finite extensions. However, there is a particularly important specialisation: $b = h^\vee$, with $(h^\vee, q) = 1$, in which case the level $k$ is \emph{boundary admissible}.\footnote{For a simply-laced Lie algebra, the level $k$ is said to be admissible if it is of the form $k = -h^\vee + \frac{p}{q}$, where $p,q \in \mathbb{Z}_{\geqslant 1}$, $p \geqslant h^\vee$ and $(p,q) = 1$. When additionally, $p = h^\vee$, the level is said to be boundary admissible.} In addition to the absence of mass parameters, the four-dimensional $D_q^{h^\vee}(\gf)$ theories in this case have no exactly marginal couplings. With the stress tensor of the VOA $\mathbb{V}[D_q^{h^\vee}(\gf)]$ being the Segal--Sugawara stress tensor of $V_{k}(\gf)$ \cite{Beem:2018duj}, certain technical niceness assumptions imply that the affine Kac--Moody subalgebra $V_{k}(\gf)$ does not admit non-trivial finite extensions \cite{Arakawa:2021ogm}. Therefore, it is also expected that all the boundary admissible affine Kac--Moody VOAs appearing in \eqref{eq:DqbgVOA} are graded unitary. Indeed, boundary admissibility has already been observed to be a recurring feature of affine Kac--Moody VOAs arising from four-dimensional unitary theories \cite{ArabiArdehali:2025fad, Beem:2026lkq}. 

Since deriving the $\mathfrak{R}$-filtration from first principles is beyond the scope of this work, we make the working assumption that the $\mathfrak{R}$-filtration on the affine Kac--Moody VOA $V_k(\gf)$ is weight-based, with appropriate adjustments for Casimir operators, as introduced in \cite{ArabiArdehali:2025fad}. It is worth emphasising that affine Kac--Moody VOAs with different, \emph{non-standard} $\mfR$-filtrations \emph{do} arise in four-dimensional unitary theories \cite{Beem:2026peg}. In all known examples, however, the putative four-dimensional theories are discrete gaugings of a parent unitary theory. In \cite{ArabiArdehali:2025fad}, it was further conjectured that the standard $\mathfrak{R}$-filtration is \emph{good}; a first contribution of this paper is to prove this conjecture for \emph{any} semisimple $\gf$. 

We then restrict to simply-laced $\gf$ and examine constraints on the level $k$ from graded unitarity. In line with \cite{ArabiArdehali:2025fad}, we utilise the closed-form expression for the vacuum determinant of affine Kac--Moody algebras \cite{gorelik2007simplicity}, and take as our diagnostic the determinant-level consequence of graded unitarity. At each conformal level $L$ (and a specific fixed $\gf$-weight), we explicitly determine the relevant structure of zeros and multiplicities of the vacuum determinant. Graded unitarity predicts a definite sign for the determinant of the Shapovalov form, determined by the $R$-charges (after refining the $\mfR$-filtration to a grading). We can compare this sign prediction with the sign implied by the closed-form expression of the vacuum determinant. 

%%%%%%%%%%%%%%%%%%%%%%%%%%%%%%%%%%%%%
\begin{table}[]
\centering
\footnotesize

\renewcommand{\arraystretch}{1.60}
\setlength{\tabcolsep}{4pt}

\setlength{\arrayrulewidth}{0.6pt}
\arrayrulecolor{black}

%\begin{tabularx}
\begin{tabularx}{\textwidth}{
|>{\centering\arraybackslash}p{0.16\textwidth}
|>{\raggedright\arraybackslash}p{0.22\textwidth}
|>{\raggedright\arraybackslash}p{0.33\textwidth}
|X|
}
\hline

\textbf{Vertex algebra} & \centering \textbf{Level $k$} & \centering \textbf{Conditions} & \textbf{SCFT realisation} \\
\hline

% A
\multirow{2}{*}{$V_k(\mathfrak{sl}_N)$} & \cellcolor{knownSCFT}$-N+\frac{N}{q}$ & \cellcolor{knownSCFT}$(N,q)=1$ & \cellcolor{knownSCFT}``class I'' $D_q^N(\mathfrak{sl}_N)$ \\
\hhline{~---}

& $-N+\frac{N}{q}$ & $(N,q)>1,\quad N\nmid q$ & \\
\hline

% D
\multirow{8}{*}{\shortstack{$V_k(\mathfrak{so}_{2N})$ \\ $(N\geqslant4)$
}} & \cellcolor{knownSCFT}$-(2N-2)+\frac{2N-2}{q}$ & \cellcolor{knownSCFT}$(2N-2,q)=1$ & \cellcolor{knownSCFT}``class I'' $D_q^{2N-2}(\mathfrak{so}_{2N})$ \\
\hhline{~---}

& \cellcolor{newSCFT} $-(2N-2)+\frac{2N-2}{q}$ & \cellcolor{newSCFT} $q\in2\mathbb N+1,\quad(2N-2,q)>1$ & \cellcolor{newSCFT} \qquad $D_q^{2N-2}(\mathfrak{so}_{2N})$ \\
\hhline{~---}

& \cellcolor{knownSCFT}$-2$ & \cellcolor{knownSCFT}$N=4$ & \cellcolor{knownSCFT}$SU(2)$ with $N_f=4$ \\
\hhline{~---}

& \cellcolor{newSCFT} $-(2N-2)+\frac{N}{q}$ & \cellcolor{newSCFT} $1\leqslant q\leqslant \frac{N-2}{2}, \ \ \frac{N}{\text{gcd}(N,q)} \in 2\mathbb{N}$ & \cellcolor{newSCFT} \qquad $D_q^{N}(\mathfrak{so}_{2N})$ \\
\hhline{~---}

& $-(2N-2)+\frac{N}{q}$ & $1\leqslant q\leqslant \frac{N-2}{2}, \ \ \frac{N}{\text{gcd}(N,q)} \in 2\mathbb{N}+1$ &  \\
\hhline{~---}

& \cellcolor{newSCFT} $-(2N-2)+\frac{2N}{3N-2}$ & \cellcolor{newSCFT} $N\in4\mathbb N$ & \cellcolor{newSCFT} \qquad $D_{\frac{3N-2}{2}}^{N}(\mathfrak{so}_{2N})$ \\
\hhline{~---}

& $-(2N-2)+\frac{2N}{3N-2}$ & $N\in4\mathbb N+2$ & \\
\hhline{~---}

& $-(2N-2)+\frac{2}{q}$ & $q\in2\mathbb N+1$, \quad $N \in 2\mathbb{N}+1$ & \\
\hhline{~---}

& $-\frac{5}{2}$ & $N=4$ & \\
\hhline{~---}

& $-\frac{22}{3}$ & $N=6$ & \\
\hline

% E6
\multirow{4}{*}{$V_k(\mathfrak e_6)$} & \cellcolor{knownSCFT}$-12+\frac{12}{q}$ & \cellcolor{knownSCFT}$(12,q)=1$ & \cellcolor{knownSCFT}``class I'' $D_q^{12}(\mathfrak e_6)$ \\
\hhline{~---}

&  \cellcolor{newSCFT} $-12+\frac{3}{q}$ & \cellcolor{newSCFT} $(3,q)=1$ &  \cellcolor{newSCFT} \qquad $D_{4q}^{12}(\mathfrak e_6)$ \\
\hhline{~---}

& \cellcolor{knownSCFT}$-3$ & \cellcolor{knownSCFT} & \cellcolor{knownSCFT}rank-one $E_6$ SCFT \\
\hhline{~---}

& \cellcolor{newSCFT} $-6$ & \cellcolor{newSCFT} & \cellcolor{newSCFT} \qquad $D_{2}^{12}(\mathfrak e_6)$ \\
\hline

% E7
\multirow{5}{*}{$V_k(\mathfrak e_7)$} & \cellcolor{knownSCFT}$-18+\frac{18}{q}$ & \cellcolor{knownSCFT}$(18,q)=1$ & \cellcolor{knownSCFT}``class I'' $D_q^{18}(\mathfrak e_7)$ \\
\hhline{~---}

& \cellcolor{newSCFT} $-18+\frac{2}{q}$ & \cellcolor{newSCFT} $(2,q)=1$ & \cellcolor{newSCFT} \qquad $D^{18}_{9q}(\mathfrak{e}_7)$ \\
\hhline{~---}

& \cellcolor{knownSCFT}$-4$ & \cellcolor{knownSCFT} & \cellcolor{knownSCFT}rank-one $E_7$ SCFT \\
\hhline{~---}

& $-8$ & & \\
\hhline{~---}

& $-11$ & & \\
\hline

% E8
\multirow{6}{*}{$V_k(\mathfrak e_8)$} & \cellcolor{knownSCFT}$-30+\frac{30}{q}$ & \cellcolor{knownSCFT}$(30,q)=1$ & \cellcolor{knownSCFT}``class I'' $D_q^{30}(\mathfrak e_8)$ \\
\hhline{~---}

& \cellcolor{knownSCFT}$-6$ & \cellcolor{knownSCFT} & \cellcolor{knownSCFT}rank-one $E_8$ SCFT \\
\hhline{~---}

& \cellcolor{newSCFT} $-10$ & \cellcolor{newSCFT} & \cellcolor{newSCFT} \qquad $D_1^{20}(\mathfrak e_8)$ \\
\hhline{~---}

& \cellcolor{newSCFT} $-18$ & \cellcolor{newSCFT} & \cellcolor{newSCFT} \qquad $D_2^{24}(\mathfrak e_8)$ \\
\hhline{~---}

& \cellcolor{newSCFT} $-22$ & \cellcolor{newSCFT} & \cellcolor{newSCFT} \qquad $D_3^{24}(\mathfrak e_8)$ \\
\hhline{~---}

& $-\frac{29}{2}$ & & \\
\hline

\end{tabularx}

\caption{Simply-laced affine Kac--Moody algebras selected by determinant-based graded unitarity constraints. Entries in pale blue are either associated VOAs of rank-one Kodaira SCFTs, or are boundary admissible affine VOAs. Entries in pale green correspond to non-admissible affine VOAs conjectured to arise in the respective $D^b_q(\gf)$ AD theories.}
\label{tab:classification}

\end{table}
%%%%%%%%%%%%%%%%%%%%%%%%%%%%%%%%%%%%%

Our main result is that this comparison leads to extremely stringent restrictions on the level $k$: under the standard $\mfR$-filtration hypothesis, we find that all levels except the discrete set of values tabulated in Table \ref{tab:classification} are incompatible with graded unitarity. This set includes the affine Kac--Moody algebras associated with the rank-one Kodaira SCFTs and the ``class I" $D_q^{h^\vee}(\gf)$ generalised Argyres--Douglas theories with simply-laced flavour symmetry. Interestingly, determinant-based graded unitarity constraints in a specific fixed $\gf$-weight sector also admit certain \emph{non-admissible} affine Kac--Moody VOAs conjecturally associated with specific $D_q^b(\gf)$ Argyres--Douglas theories. Although finer state-based graded unitarity constraints may rule out some of these levels, their survival provides suggestive evidence for the respective conjectured identifications in \eqref{eq:DqbgVOA}. Notably, all the $D_q^b(\gf)$ Argyres--Douglas theories making an appearance in Table \ref{tab:classification} have no additional mass parameters/flavour symmetries arising from the irregular puncture.

More broadly, the present results reinforce an emerging picture that four-dimensional unitarity serves as a remarkably powerful rigidity principle, sharply constraining the vertex algebras appearing in the image of the SCFT/VOA map. In all the examples studied so far, Virasoro, $\mathcal{W}_3$, and now all simply-laced affine VOAs, even the determinant-based consequences of graded unitarity already single out a limited, discrete set of values that can possibly arise from four-dimensional theories. We expect similar classifying results to hold for the non-simply laced affine VOAs as well. These results, while capturing a coarse shadow of the full positivity constraints encoded by graded unitarity, point towards the need for an intrinsic vertex-algebraic formulation of graded unitarity. A central question is to determine which structural properties of a vertex algebra govern the existence of a compatible $\mfR$-filtration and conjugation satisfying the appropriate positivity conditions. 

The organisation of the paper is as follows. In Section \ref{sec:GKvacdet} we review the relevant aspects of the Gorelik--Kac vacuum determinant for simply-laced affine Kac--Moody algebras, and reformulate the functions encoding the multiplicities of zeros in a form better suited for explicit computations. In Section \ref{sec:grunitstructures} we specialise the general framework of graded unitarity to affine Kac--Moody VOAs, describe the standard $\mfR$-filtration, and prove its goodness. We also state the determinant-based sign constraint in the simply laced case. In Section \ref{sec:typeAgrunit} we analyse graded unitarity constraints for $V_k(\mathfrak{sl}_N)$ affine Kac--Moody VOAs. We begin with a bespoke low-level analysis, and then study the structure of zeros of the vacuum determinant to deduce the constraints on the level $k$ and obtain the aforementioned restriction. Section \ref{sec:typeDgrunit} carries out the analogous analysis for $V_k(\mathfrak{so}_{2N})$. After treating $V_k(\mathfrak{so}_8)$ explicitly, we analyse the case of even $N$, and then summarise the complete classification, with the details for odd $N$ relegated to Appendix \ref{app:VkSOodd}. Finally, in Section \ref{sec:typeEgrunit} we complete the characterisation of graded unitarity constraints for simply-laced affine Kac--Moody VOAs by determining the determinant-based constraints for $V_k(\mathfrak{e}_{n=6,7,8})$. 